\documentclass{amsart}
% For dealing with references we use the comment environment
\usepackage{verbatim}
\newenvironment{reference}{\comment}{\endcomment}
%\newenvironment{reference}{}{}
\newenvironment{slogan}{\comment}{\endcomment}
\newenvironment{history}{\comment}{\endcomment}

% For commutative diagrams we use Xy-pic
\usepackage[all]{xy}

% We use 2cell for 2-commutative diagrams.
\xyoption{2cell}
\UseAllTwocells

% We use multicol for the list of chapters between chapters
\usepackage{multicol}

% This is generally recommended for better output
\usepackage{lmodern}
\usepackage[T1]{fontenc}

% For cross-file-references

% Package for hypertext links:
\usepackage{hyperref}

% For any local file, say "hello.tex" you want to link to please

% Theorem environments.
%
\theoremstyle{plain}
\newtheorem{theorem}[subsection]{Theorem}
\newtheorem{proposition}[subsection]{Proposition}
\newtheorem{lemma}[subsection]{Lemma}

\theoremstyle{definition}
\newtheorem{definition}[subsection]{Definition}
\newtheorem{example}[subsection]{Example}
\newtheorem{exercise}[subsection]{Exercise}
\newtheorem{situation}[subsection]{Situation}

\theoremstyle{remark}
\newtheorem{remark}[subsection]{Remark}
\newtheorem{remarks}[subsection]{Remarks}

\numberwithin{equation}{subsection}

% Macros
%
\def\lim{\mathop{\mathrm{lim}}\nolimits}
\def\colim{\mathop{\mathrm{colim}}\nolimits}
\def\Spec{\mathop{\mathrm{Spec}}}
\def\Hom{\mathop{\mathrm{Hom}}\nolimits}
\def\Ext{\mathop{\mathrm{Ext}}\nolimits}
\def\SheafHom{\mathop{\mathcal{H}\!\mathit{om}}\nolimits}
\def\SheafExt{\mathop{\mathcal{E}\!\mathit{xt}}\nolimits}
\def\Sch{\mathit{Sch}}
\def\Mor{\mathop{\mathrm{Mor}}\nolimits}
\def\Ob{\mathop{\mathrm{Ob}}\nolimits}
\def\Sh{\mathop{\mathit{Sh}}\nolimits}
\def\NL{\mathop{N\!L}\nolimits}
\def\CH{\mathop{\mathrm{CH}}\nolimits}
\def\proetale{{pro\text{-}\acute{e}tale}}
\def\etale{{\acute{e}tale}}
\def\QCoh{\mathit{QCoh}}
\def\Ker{\mathop{\mathrm{Ker}}}
\def\Im{\mathop{\mathrm{Im}}}
\def\Coker{\mathop{\mathrm{Coker}}}
\def\Coim{\mathop{\mathrm{Coim}}}

% Boxtimes
%
\DeclareMathSymbol{\boxtimes}{\mathbin}{AMSa}{"02}

%
% Macros for moduli stacks/spaces
%
\def\QCohstack{\mathcal{QC}\!\mathit{oh}}
\def\Cohstack{\mathcal{C}\!\mathit{oh}}
\def\Spacesstack{\mathcal{S}\!\mathit{paces}}
\def\Quotfunctor{\mathrm{Quot}}
\def\Hilbfunctor{\mathrm{Hilb}}
\def\Curvesstack{\mathcal{C}\!\mathit{urves}}
\def\Polarizedstack{\mathcal{P}\!\mathit{olarized}}
\def\Complexesstack{\mathcal{C}\!\mathit{omplexes}}
% \Pic is the operator that assigns to X its picard group, usage \Pic(X)
% \Picardstack_{X/B} denotes the Picard stack of X over B
% \Picardfunctor_{X/B} denotes the Picard functor of X over B
\def\Pic{\mathop{\mathrm{Pic}}\nolimits}
\def\Picardstack{\mathcal{P}\!\mathit{ic}}
\def\Picardfunctor{\mathrm{Pic}}
\def\Deformationcategory{\mathcal{D}\!\mathit{ef}}

\setlength{\emergencystretch}{3em}
\begin{document}
\title[Stacks Project, Tag 0EFC]{310 --- Faltings annihilator theorem for local cohomology}
\maketitle
\noindent Copyright (C) 2005--2025 Johan de Jong. Stacks Project authors. Source: \href{https://stacks.math.columbia.edu/tag/0EFC}{Tag 0EFC}.

\noindent Editorial difficulty note (not author text). Difficulty H3: The proof constructs a global annihilator by a double induction on cohomological degree and dimension-function strata. Spectral-sequence filtrations control successive powers of annihilators, and finite dualizing Ext supports produce uniform choices. The converse locates a minimal failing depth-codimension pair and proves its local cohomology cannot be annihilated..

\noindent Editorial provenance note (not author text). History: excerpt from revision \texttt{a0444\allowbreak 6e57e\allowbreak c1fbc\allowbreak 252a8\allowbreak 71afc\allowbreak ec775\allowbreak 2fb28\allowbreak 07b14}, local-cohomology.tex, lines 2034--2206. Reference presentation was adapted; original-author context and core support blocks were retained or added. The mathematical wording is unchanged. Licensed under GNU FDL 1.2 or later; no invariant sections or cover texts. See ../licenses/COPYING. Context blocks reproduce author definitions and supporting results; they are not additional counted problems.

\section{Dualizing complexes and dimension functions}
\label{dualizing-section-dimension-function}

\noindent
Our results in the local setting have the following consequence:
a Noetherian ring which has a dualizing complex is a
universally catenary ring of finite dimension.



\begin{proposition}
\label{proposition-annihilator}
\begin{reference}
\href{https://stacks.math.columbia.edu/bibliography}{Bibliography: Faltings-annulators}.
\end{reference}
Let $A$ be a Noetherian ring which has a dualizing complex.
Let $T \subset T' \subset \Spec(A)$ be subsets stable under
specialization. Let $s \geq 0$ an integer. Let $M$ be a finite $A$-module.
The following are equivalent
\begin{enumerate}
\item there exists an ideal $J \subset A$ with $V(J) \subset T'$
such that $J$ annihilates $H^i_T(M)$ for $i \leq s$, and
\item for all $\mathfrak p \not \in T'$,
$\mathfrak q \in T$ with $\mathfrak p \subset \mathfrak q$
we have
$$
\text{depth}_{A_\mathfrak p}(M_\mathfrak p) +
\dim((A/\mathfrak p)_\mathfrak q) > s
$$
\end{enumerate}
\end{proposition}

\begin{proof}
Let $\omega_A^\bullet$ be a dualizing complex. Let $\delta$ be its
dimension function, see Dualizing Complexes, Section
\ref{dualizing-section-dimension-function}.
An important role will be played by the finite $A$-modules
$$
E^i = \Ext_A^i(M, \omega_A^\bullet)
$$
For $\mathfrak p \subset A$ we will write $H^i_\mathfrak p$ to denote the
local cohomology of an $A_\mathfrak p$-module with respect to
$\mathfrak pA_\mathfrak p$. Then we see that
the $\mathfrak pA_\mathfrak p$-adic completion of
$$
(E^i)_\mathfrak p =
\Ext^{\delta(\mathfrak p) + i}_{A_\mathfrak p}(M_\mathfrak p,
(\omega_A^\bullet)_\mathfrak p[-\delta(\mathfrak p)])
$$
is Matlis dual to
$$
H^{-\delta(\mathfrak p) - i}_{\mathfrak p}(M_\mathfrak p)
$$
by
Dualizing Complexes, Lemma \href{https://stacks.math.columbia.edu/tag/0AAK}{lemma special case local duality (Tag 0AAK)}.
In particular we deduce from this the
following fact: an ideal $J \subset A$ annihilates
$(E^i)_\mathfrak p$ if and only if $J$ annihilates
$H^{-\delta(\mathfrak p) - i}_{\mathfrak p}(M_\mathfrak p)$.

\medskip\noindent
Set $T_n = \{\mathfrak p \in T \mid \delta(\mathfrak p) \leq n\}$.
As $\delta$ is a bounded function, we see that
$T_a = \emptyset$ for $a \ll 0$ and $T_b = T$ for $b \gg 0$.

\medskip\noindent
Assume (2). Let us prove the existence of $J$ as in (1).
We will use a double induction to do this. For $i \leq s$
consider the induction hypothesis $IH_i$:
$H^a_T(M)$ is annihilated by some $J \subset A$ with $V(J) \subset T'$
for $0 \leq a \leq i$. The case $IH_0$ is trivial
because $H^0_T(M)$ is a submodule of $M$ and hence finite
and hence is annihilated by some ideal $J$ with $V(J) \subset T$.

\medskip\noindent
Induction step. Assume $IH_{i - 1}$ holds for some $0 < i \leq s$.
Pick $J'$ with $V(J') \subset T'$ annihilating $H^a_T(M)$ for
$0 \leq a \leq i - 1$ (the induction hypothesis guarantees we can
do this). We will show by descending induction on $n$
that there exists an ideal $J$ with $V(J) \subset T'$ such that the
associated primes of $J H^i_T(M)$ are in $T_n$.
For $n \ll 0$ this implies $JH^i_T(M) = 0$ 
(Algebra, Lemma \href{https://stacks.math.columbia.edu/tag/0587}{lemma ass zero (Tag 0587)})
and hence $IH_i$ will hold.
The base case $n \gg 0$ is trivial because $T = T_n$ in this case
and all associated primes of $H^i_T(M)$ are in $T$.

\medskip\noindent
Thus we assume given $J$ with the property for $n$.
Let $\mathfrak q \in T_n$. Let $T_\mathfrak q \subset \Spec(A_\mathfrak q)$
be the inverse image of $T$. We have
$H^j_T(M)_\mathfrak q = H^j_{T_\mathfrak q}(M_\mathfrak q)$
by Lemma \href{https://stacks.math.columbia.edu/tag/0EF5}{lemma torsion change rings (Tag 0EF5)}.
Consider the spectral sequence
$$
H_\mathfrak q^p(H^q_{T_\mathfrak q}(M_\mathfrak q))
\Rightarrow
H^{p + q}_\mathfrak q(M_\mathfrak q)
$$
of Lemma \href{https://stacks.math.columbia.edu/tag/0EF6}{lemma local cohomology ss (Tag 0EF6)}.
Below we will find an ideal $J'' \subset A$ with $V(J'') \subset T'$
such that $H^i_\mathfrak q(M_\mathfrak q)$ is annihilated by $J''$ for all
$\mathfrak q \in T_n \setminus T_{n - 1}$.
Claim: $J (J')^i J''$ will work for $n - 1$.
Namely, let $\mathfrak q \in T_n \setminus T_{n - 1}$.
The spectral sequence above defines a filtration
$$
E_\infty^{0, i} = E_{i + 2}^{0, i} \subset \ldots \subset E_3^{0, i} \subset
E_2^{0, i} = H^0_\mathfrak q(H^i_{T_\mathfrak q}(M_\mathfrak q))
$$
The module $E_\infty^{0, i}$ is annihilated by $J''$.
The subquotients $E_j^{0, i}/E_{j + 1}^{0, i}$ for $i + 1 \geq j \geq 2$
are annihilated by $J'$ because the target of $d_j^{0, i}$
is a subquotient of
$$
H^j_\mathfrak q(H^{i - j + 1}_{T_\mathfrak q}(M_\mathfrak q)) =
H^j_\mathfrak q(H^{i - j + 1}_T(M)_\mathfrak q)
$$
and $H^{i - j + 1}_T(M)_\mathfrak q$ is annihilated by $J'$ by choice of $J'$.
Finally, by our choice of $J$ we have
$J H^i_T(M)_\mathfrak q \subset H^0_\mathfrak q(H^i_T(M)_\mathfrak q)$
since the non-closed points of $\Spec(A_\mathfrak q)$ have higher
$\delta$ values. Thus $\mathfrak q$ cannot be an associated prime of
$J(J')^iJ'' H^i_T(M)$ as desired.

\medskip\noindent
By our initial remarks we see that $J''$ should annihilate
$$
(E^{-\delta(\mathfrak q) - i})_\mathfrak q =
(E^{-n - i})_\mathfrak q
$$
for all $\mathfrak q \in T_n \setminus T_{n - 1}$.
But if $J''$ works for one $\mathfrak q$, then it works for all
$\mathfrak q$ in an open neighbourhood of $\mathfrak q$
as the modules $E^{-n - i}$ are finite.
Since every subset of $\Spec(A)$ is Noetherian with the induced
topology (Topology, Lemma \href{https://stacks.math.columbia.edu/tag/0052}{lemma Noetherian (Tag 0052)}),
we conclude that it suffices
to prove the existence of $J''$ for one $\mathfrak q$.

\medskip\noindent
Since the ext modules are finite the existence of $J''$ is
equivalent to
$$
\text{Supp}(E^{-n - i}) \cap \Spec(A_\mathfrak q) \subset T'.
$$
This is equivalent to showing the localization of $E^{-n - i}$ at every
$\mathfrak p \subset \mathfrak q$, $\mathfrak p \not \in T'$
is zero. Using local duality over $A_\mathfrak p$ we find that we need
to prove that
$$
H^{i + n - \delta(\mathfrak p)}_\mathfrak p(M_\mathfrak p) =
H^{i - \dim((A/\mathfrak p)_\mathfrak q)}_\mathfrak p(M_\mathfrak p)
$$
is zero (this uses that $\delta$ is a dimension function).
This vanishes by the assumption in the lemma and $i \leq s$ and
Dualizing Complexes, Lemma \href{https://stacks.math.columbia.edu/tag/0AVZ}{lemma depth (Tag 0AVZ)}.

\medskip\noindent
To prove the converse implication we assume (2) does not hold
and we work backwards through the arguments above. First, we pick a
$\mathfrak q \in T$, $\mathfrak p \subset \mathfrak q$
with $\mathfrak p \not \in T'$ such that
$$
i = \text{depth}_{A_\mathfrak p}(M_\mathfrak p) +
\dim((A/\mathfrak p)_\mathfrak q) \leq s
$$
is minimal. Then
$H^{i - \dim((A/\mathfrak p)_\mathfrak q)}_\mathfrak p(M_\mathfrak p)$
is nonzero by the nonvanishing in
Dualizing Complexes, Lemma \href{https://stacks.math.columbia.edu/tag/0AVZ}{lemma depth (Tag 0AVZ)}.
Set $n = \delta(\mathfrak q)$. Then
there does not exist an ideal $J \subset A$ with $V(J) \subset T'$
such that $J(E^{-n - i})_\mathfrak q = 0$.
Thus $H^i_\mathfrak q(M_\mathfrak q)$ is not
annihilated by an ideal $J \subset A$ with $V(J) \subset T'$.
By minimality of $i$ it follows from the spectral sequence displayed above
that the module $H^i_T(M)_\mathfrak q$
is not annihilated by an ideal $J \subset A$
with $V(J) \subset T'$. Thus $H^i_T(M)$
is not annihilated by an ideal $J \subset A$
with $V(J) \subset T'$. This finishes the proof of the proposition.
\end{proof}
\end{document}
